\documentclass[aspectratio=169,11pt]{beamer}

\usetheme{Madrid}
\usecolortheme{default}
\usefonttheme{professionalfonts}
\setbeamertemplate{navigation symbols}{}
\setbeamertemplate{footline}[frame number]

\usepackage{amsmath, amssymb, booktabs}

\title{Health, Risk, And Labor Markets}
\subtitle{Core Frameworks And Measurement}
\author{Special Topic 6: Labor Market, Health, and Population -- Week 1}
\date{}

\begin{document}

\begin{frame}
  \titlepage
\end{frame}

\begin{frame}{Central question}
\begin{itemize}
    \item How should labor economists model and measure health when health affects both productivity and the feasible set for work?
    \item Health is not only a control variable. It changes work capacity, job choice, productivity, mobility, wages, amenities, and welfare.
    \item Week 1 discipline: name the health object before interpreting a health coefficient.
\end{itemize}
\end{frame}

\begin{frame}{Why this is labor economics}
\begin{itemize}
    \item Health changes whether work is feasible and how much work is sustainable.
    \item Health changes which jobs, schedules, commutes, risks, and amenities are tolerable.
    \item Health changes productivity inside jobs and firm decisions about hiring, tasks, scheduling, insurance, and accommodation.
    \item Labor outcomes are the organizing objects: employment, hours, job choice, wages, productivity, mobility, firm behavior, and welfare.
\end{itemize}
\end{frame}

\begin{frame}{Health as human capital and work capacity}
\small
\begin{tabular}{p{0.24\linewidth} p{0.42\linewidth} p{0.22\linewidth}}
\toprule
Dimension & Labor interpretation & Typical proxy \\
\midrule
Work capacity & ability to sustain hours and tasks & limitations, disability, function \\
Productivity & output conditional on working & performance, absence, demands \\
Risk & future impairment, exit, mortality & shocks, conditions, mortality risk \\
Feasible set & tolerable jobs and schedules & occupation, hours, mobility \\
Welfare & value of work net of pain, stress, risk & amenities, insurance, job quality \\
\bottomrule
\end{tabular}
\vspace{0.3cm}
\begin{itemize}
    \item Health shifts both the budget set and the feasible set.
\end{itemize}
\end{frame}

\begin{frame}{Measurement choices}
\small
\begin{tabular}{p{0.28\linewidth} p{0.38\linewidth} p{0.22\linewidth}}
\toprule
Measure & What it may capture & Main warning \\
\midrule
Self-reported health & pain, fatigue, private functioning & reporting and justification bias \\
Diagnosis / biomarker & clinical status or severity & may miss work capacity \\
Disability status & health plus administrative criteria & blends health and program rules \\
Hospitalization / shock & severe discrete event & narrow margin \\
Mental health scale & cognition, stress, persistence & latent and stigmatized \\
Insurance / treatment & access and recovery margin & jointly affects work and health \\
\bottomrule
\end{tabular}
\end{frame}

\begin{frame}{Different measures, different objects}
\begin{itemize}
    \item Self-reports can reveal private information about function, pain, fatigue, or latent conditions.
    \item Diagnoses and biomarkers can look objective but may miss severity, untreated illness, and task feasibility.
    \item Disability status is an institutional outcome, not pure work capacity.
    \item Mortality risk and acute shocks identify severe events, not routine chronic constraints.
    \item Insurance and treatment access can change health, reported health, labor supply, and mobility at the same time.
\end{itemize}
\end{frame}

\begin{frame}{Reverse causality and dynamic selection}
\small
\[
Y_{it} =
\beta H_{it}
+ X_{it}'\Gamma
+ \alpha_i
+ \lambda_t
+ \varepsilon_{it}
\]
\begin{itemize}
    \item Reverse causality: work affects health through stress, injuries, income, insurance, schedules, and job conditions.
    \item Omitted SES: education, wealth, early-life resources, occupation, and local labor demand shape both health and work.
    \item Dynamic selection: healthier workers remain employed longer and sort into different firms, jobs, and schedules.
    \item Program rules: disability and insurance institutions can move labor supply independently of latent health.
\end{itemize}
\end{frame}

\begin{frame}{Worker-side versus firm-side channels}
\small
\begin{tabular}{p{0.28\linewidth} p{0.31\linewidth} p{0.28\linewidth}}
\toprule
Health channel & Worker-side margin & Firm-side margin \\
\midrule
Reduced capacity & hours, exit, job downgrading & screening, reassignment, accommodation \\
Higher risk & safer jobs, retirement timing & insurance, safety investment \\
Treatment access & recovery, job lock, mobility & benefit design, job packages \\
Uncertainty & precautionary supply and saving & absenteeism and scheduling risk \\
Chronic burden & commuting and mobility limits & flexibility, monitoring, retention \\
\bottomrule
\end{tabular}
\vspace{0.25cm}
\begin{itemize}
    \item Reduced-form effects can mix supply, productivity, sorting, and employer response.
\end{itemize}
\end{frame}

\begin{frame}{A field map for the course}
\small
\begin{tabular}{p{0.30\linewidth} p{0.52\linewidth}}
\toprule
Later object & Week 1 preparation \\
\midrule
Disability and chronic conditions & separates health from program status and names the labor margin \\
Fertility, mortality, caregiving & treats shocks as lifecycle labor-allocation events \\
Mental health and productivity & makes latent conditions and welfare measurement explicit \\
Disease and environment & scales health from individual heterogeneity to market condition \\
Research synthesis & uses labor object -> health object -> design problem -> welfare object \\
\bottomrule
\end{tabular}
\end{frame}

\begin{frame}{Research Lab design}
\begin{itemize}
    \item Primary anchor: Blundell, Britton, Costa Dias, and French on health measures and employment near retirement.
    \item Challenge anchor: Bound on self-reported versus objective health measures.
    \item Reproduce: compare reduced-form labor gradients across health measures.
    \item Diagnose: classify measurement objects and threats from reverse causality, SES, and selection.
    \item Transfer: move the logic to acute shocks, disability onset, mental health treatment, or insurance access.
\end{itemize}
\end{frame}

\begin{frame}{Research design checklist}
\begin{enumerate}
    \item Name the labor margin: employment, hours, wages, productivity, job choice, mobility, amenities, or welfare.
    \item Name the health object: self-report, diagnosis, biomarker, disability status, severe shock, mental health, insurance, or treatment.
    \item State the timing and selection problem.
    \item Separate worker-side and firm-side channels.
    \item Say what welfare object is visible and what remains offstage.
\end{enumerate}
\end{frame}

\begin{frame}{Bridge to Week 2}
\begin{itemize}
    \item Week 1 builds the common language: health as capacity, risk, feasible set, measurement problem, and labor-market margin.
    \item Week 2 applies that language to disability, chronic conditions, job choice, accommodations, and program interaction.
    \item The next question: when does a health limitation become an employment constraint, a job-quality constraint, or an institutional response?
\end{itemize}
\end{frame}

\begin{frame}{Takeaways}
\begin{itemize}
    \item Health is a labor object when it changes capacity, productivity, risk, feasible work, firm behavior, or welfare.
    \item Measurement is not a technical afterthought; it is one of the central open problems in the field.
    \item Reverse causality, omitted SES, and dynamic selection are baseline threats, not footnotes.
    \item Strong designs connect a health object to a labor margin, a channel, a counterfactual, and a welfare interpretation.
\end{itemize}
\end{frame}

\end{document}
